\documentclass[12pt,a4paper]{article}

\usepackage[utf8]{inputenc}
\usepackage[T1]{fontenc}
\usepackage[brazil]{babel}
\usepackage[top=2cm, bottom=2cm, left=2cm, right=2.5cm]{geometry}
\usepackage{amsmath}
\usepackage{amssymb}
\usepackage{xcolor}
\usepackage{tcolorbox}
\tcbuselibrary{skins, breakable}
\usepackage{enumitem}
\usepackage{booktabs}
\usepackage{array}
\usepackage{tabularx}
\usepackage{multirow}
\usepackage{hyperref}
\hypersetup{colorlinks=true, linkcolor=azulphysics}
\usepackage{fontawesome5}
\usepackage{tikz}
\usepackage{pgfplots}
\pgfplotsset{compat=1.18}

% ---------------------------------------------------------------
% Cores
% ---------------------------------------------------------------
\definecolor{azulphysics}{RGB}{0, 82, 155}
\definecolor{verdeok}{RGB}{0, 128, 64}
\definecolor{laranjaalerta}{RGB}{200, 100, 0}
\definecolor{vermelhocorte}{RGB}{180, 30, 30}
\definecolor{roxoprojetor}{RGB}{100, 0, 160}
\definecolor{cinzafundo}{RGB}{245, 245, 245}

% ---------------------------------------------------------------
% Boxes
% ---------------------------------------------------------------
\tcbset{
  blocobox/.style={
    colback=azulphysics!8!white,
    colframe=azulphysics,
    fonttitle=\bfseries\large,
    title={#1},
    arc=3mm,
    left=6pt, right=6pt, top=6pt, bottom=6pt,
    breakable
  },
  quadrobox/.style={
    colback=white,
    colframe=black!60,
    fonttitle=\bfseries\small,
    title={\faChalkboard\ Quadro branco},
    arc=2mm,
    left=5pt, right=5pt, top=4pt, bottom=4pt,
    breakable
  },
  projetorbox/.style={
    colback=roxoprojetor!6!white,
    colframe=roxoprojetor!70!black,
    fonttitle=\bfseries\small,
    title={\faProjectDiagram\ Projetor},
    arc=2mm,
    left=5pt, right=5pt, top=4pt, bottom=4pt,
    breakable
  },
  perguntabox/.style={
    colback=verdeok!8!white,
    colframe=verdeok!70!black,
    fonttitle=\bfseries\small,
    title={\faComments\ Perguntas para a turma},
    arc=2mm,
    left=5pt, right=5pt, top=4pt, bottom=4pt
  },
  dicabox/.style={
    colback=laranjaalerta!8!white,
    colframe=laranjaalerta!70!black,
    fonttitle=\bfseries\small,
    title={\faLightbulb\ Dica did\'atica},
    arc=2mm,
    left=5pt, right=5pt, top=4pt, bottom=4pt
  },
  cortebox/.style={
    colback=vermelhocorte!6!white,
    colframe=vermelhocorte!70!black,
    fonttitle=\bfseries\small,
    title={\faClock\ Se o tempo apertar},
    arc=2mm,
    left=5pt, right=5pt, top=4pt, bottom=4pt
  }
}

% Marcador de tempo
\newcommand{\tempo}[1]{%
  \begin{flushright}
    \colorbox{azulphysics!20}{\small\textbf{\faClock\ \textasciitilde{}#1 min}}
  \end{flushright}\vspace{-0.5em}
}

% ---------------------------------------------------------------
\title{
  \textbf{\large Cursinho Solid\'ario --- F\'isica}\\[0.5em]
  \textbf{\Huge Plano de Aula}\\[0.3em]
  \textbf{\Large Cinem\'atica}\\[0.5em]
  \normalsize \textit{Guia do Professor}
}
\author{}
\date{2026}

% ---------------------------------------------------------------
\begin{document}
% ---------------------------------------------------------------

\maketitle
\thispagestyle{empty}

% ---------------------------------------------------------------
% VISÃO GERAL
% ---------------------------------------------------------------
\begin{tcolorbox}[
  colback=cinzafundo, colframe=azulphysics!50,
  arc=3mm, left=8pt, right=8pt, top=8pt, bottom=8pt
]
\begin{center}
  \textbf{\large Vis\~ao Geral da Aula}
\end{center}

\renewcommand{\arraystretch}{1.4}
\begin{tabularx}{\linewidth}{lX}
  \toprule
  \textbf{Disciplina} & F\'isica --- Cinem\'atica \\
  \textbf{P\'ublico} & Vestibulandos (ENEM e UTFPR) \\
  \textbf{Formato} & Expositivo dialogado no quadro + projetor para figuras \\
  \textbf{Estrutura} & 6 blocos modulares independentes \\
  \midrule
  \textbf{Tempo total (todos os blocos)} & $\approx 3$h \\
  \textbf{Vers\~ao reduzida (blocos obrigat\'orios)} & $\approx 1$h45 \\
  \midrule
  \multicolumn{2}{l}{\textbf{Materiais}} \\
  \multicolumn{2}{l}{%
    \begin{minipage}{\linewidth}
    \begin{itemize}[nosep, leftmargin=1.5em]
      \item Quadro branco e marcadores (pelo menos 2 cores)
      \item Projetor + notebook com o arquivo
        \texttt{notas\_aula\_cinematica.pdf}
      \item Apostila impressa para os alunos (opcional)
    \end{itemize}
    \end{minipage}
  } \\
  \bottomrule
\end{tabularx}

\medskip
\textbf{Como usar este plano:}
Cada bloco indica o tempo estimado no canto superior direito.
Os blocos marcados com \colorbox{verdeok!30}{\small ESSENCIAL} cobrem o m\'inimo
necess\'ario para ENEM. Os demais s\~ao importantes mas cortáveis se o
hor\'ario apertar.
\end{tcolorbox}

\vspace{0.8em}

% ---------------------------------------------------------------
% CRONOGRAMA RESUMO
% ---------------------------------------------------------------
\begin{tcolorbox}[colback=white, colframe=azulphysics!40, arc=2mm,
  fonttitle=\bfseries, title={Cronograma Resumido}]

\renewcommand{\arraystretch}{1.5}
\begin{tabularx}{\linewidth}{clXc}
  \toprule
  \textbf{Bloco} & \textbf{T\'opico} & \textbf{Foco principal} & \textbf{Tempo} \\
  \midrule
  0 & Abertura & Nivelamento + contrato did\'atico
    & 10--15 min \\
  1 & Conceitos fundamentais & Referencial, deslocamento, $\bar{v}$, $\bar{a}$
    & 20--25 min \\
  2 & MRU & Equa\c{c}\~ao hor\'aria, gr\'aficos, exemplo de encontro
    & 25--35 min \\
  3 & MRUV & Dedu\c{c}\~ao das 3 equa\c{c}\~oes, gr\'aficos, frenagem
    & 40--50 min \\
  4 & Queda livre e lan\c{c}. vertical & Galileu, equa\c{c}\~oes, simetria
    & 20--25 min \\
  5 & Lan\c{c}amento obl\'iquo & Decomposi\c{c}\~ao, R, H, T
    & 35--40 min \\
  6 & Fixação com ENEM & 3--4 quest\~oes comentadas
    & 20--25 min \\
  \midrule
  & \textbf{Total} & & \textbf{2h50--3h15} \\
  \bottomrule
\end{tabularx}
\end{tcolorbox}

\newpage
\tableofcontents
\newpage

% ===================================================================
\section*{Bloco 0 \hfill\colorbox{verdeok!30}{\small ESSENCIAL}}
\addcontentsline{toc}{section}{Bloco 0 --- Abertura e Nivelamento}
\section*{\hspace{1.2em}Abertura e Nivelamento}
% ===================================================================
\tempo{10--15}

\begin{tcolorbox}[blocobox={Objetivo}]
  Quebrar o gelo, identificar o n\'ivel da turma e estabelecer como a aula
  vai funcionar (quadro + projetor, participa\c{c}\~ao ativa).
\end{tcolorbox}

\begin{tcolorbox}[perguntabox]
  Lance estas perguntas oralmente e deixe os alunos responderem antes de
  formalizar qualquer coisa:
  \begin{enumerate}[nosep]
    \item ``O que \'e velocidade para voc\^es?''
    \item ``Como voc\^es saberiam dizer se um objeto est\'a em movimento?''
    \item ``Algu\'em sabe o que \'e acelera\c{c}\~ao? D\^e um exemplo do dia a dia.''
  \end{enumerate}
  \textit{N\~ao corrija erros agora. Apenas ouça e anote no quadro as ideias
  que surgem --- voc\^e vai retomar no Bloco 1.}
\end{tcolorbox}

\begin{tcolorbox}[quadrobox]
  Escreva as ideias que os alunos trouxeram nas perguntas acima (mesmo as
  erradas). Isso cria um ``painel de ideias pr\'evias'' que voc\^e vai
  refinando ao longo da aula.
\end{tcolorbox}

\begin{tcolorbox}[dicabox]
  \textbf{Contrato did\'atico:} Explique que a aula ser\'a no quadro.
  Vocês vão deduzir as fórmulas juntos --- não é para copiar antes de
  entender. Encoraje perguntas a qualquer momento.
\end{tcolorbox}

\newpage

% ===================================================================
\section*{Bloco 1 \hfill\colorbox{verdeok!30}{\small ESSENCIAL}}
\addcontentsline{toc}{section}{Bloco 1 --- Conceitos Fundamentais}
\section*{\hspace{1.2em}Conceitos Fundamentais}
% ===================================================================
\tempo{20--25}

\begin{tcolorbox}[blocobox={Objetivos}]
  Alunos devem sair sabendo distinguir:
  deslocamento vs.\ dist\^ancia percorrida,
  velocidade m\'edia vs.\ velocidade instant\^anea,
  e o papel do referencial.
\end{tcolorbox}

\begin{tcolorbox}[quadrobox]
  \textbf{1. Referencial e posi\c{c}\~ao} (3--4 min)
  \begin{itemize}[nosep]
    \item Desenhe um eixo horizontal no quadro. Marque um ponto O
          (referencial) e uma posi\c{c}\~ao $x$.
    \item Escreva: \textbf{posi\c{c}\~ao} $= x$; positivo \`a direita,
          negativo \`a esquerda.
  \end{itemize}

  \medskip
  \textbf{2. Deslocamento vs.\ dist\^ancia} (5 min)
  \begin{itemize}[nosep]
    \item Desenhe: carro vai de $x_i = 2$ m at\'e $x_f = 8$ m e volta para
          $x = 5$ m.
    \item Quadro: $\Delta x = x_f - x_i = 5 - 2 = 3$ m
          (deslocamento $\neq$ dist\^ancia percorrida $= 9$ m).
    \item Conclua: \textbf{deslocamento} tem sinal; \textbf{dist\^ancia}
          \'e sempre positiva.
  \end{itemize}

  \medskip
  \textbf{3. Velocidade m\'edia e acelera\c{c}\~ao m\'edia} (5--7 min)
  \begin{itemize}[nosep]
    \item Escreva as defini\c{c}\~oes lado a lado:
    \[
      \bar{v} = \frac{\Delta x}{\Delta t} \qquad
      \bar{a} = \frac{\Delta v}{\Delta t}
    \]
    \item Destaque a convers\~ao: $1$ m/s $= 3{,}6$ km/h.
    \item Exemplo r\'apido: carro percorre 200 m em 10 s $\to$
          $\bar{v} = 20$ m/s $= 72$ km/h.
  \end{itemize}
\end{tcolorbox}

\begin{tcolorbox}[projetorbox]
  \textbf{Figura 1:} Eixo com referencial, posi\c{c}\~oes positivas e
  negativas, e seta de deslocamento.
  Use a p\'agina 3 das notas de aula (\texttt{notas\_aula\_cinematica.pdf}).

  \medskip
  Projete \textbf{apenas quando precisar que os alunos visualizem o eixo
  com clareza} --- prefira desenhar no quadro sempre que possível.
\end{tcolorbox}

\begin{tcolorbox}[perguntabox]
  ``Se um carro vai do km 10 ao km 30 e volta ao km 20 da rodovia,
  qual o deslocamento? Qual a dist\^ancia percorrida?''

  \textit{Resposta esperada: $\Delta x = 10$ km; $d = 30$ km.
  Corrija quem confundir os dois conceitos.}
\end{tcolorbox}

\begin{tcolorbox}[dicabox]
  O erro mais comum nos primeiros contatos com cinem\'atica \'e usar
  dist\^ancia no lugar de deslocamento (e vice-versa). Frise isso agora
  para evitar erros em todo o resto da aula.
\end{tcolorbox}

\begin{tcolorbox}[cortebox]
  Se precisar de tempo: suprima o exemplo num\'erico de $\bar{v}$ ---
  ele vai reaparecer no Bloco 2.
\end{tcolorbox}

\newpage

% ===================================================================
\section*{Bloco 2 \hfill\colorbox{verdeok!30}{\small ESSENCIAL}}
\addcontentsline{toc}{section}{Bloco 2 --- MRU}
\section*{\hspace{1.2em}MRU --- Movimento Retil\'ineo Uniforme}
% ===================================================================
\tempo{25--35}

\begin{tcolorbox}[blocobox={Objetivos}]
  Deduzir a equação horária a partir da definição de $\bar v$.
  Ler e interpretar gráficos $x \times t$ e $v \times t$.
  Resolver o problema clássico de encontro de dois móveis.
\end{tcolorbox}

\begin{tcolorbox}[quadrobox]
  \textbf{1. Defini\c{c}\~ao e dedu\c{c}\~ao} (5--7 min)
  \begin{itemize}[nosep]
    \item Escreva no quadro: ``No MRU, $a = 0$ e $v = $ cte.''
    \item Dedu\c{c}\~ao ao vivo (escreva cada passo):
    \[
      v = \frac{x - x_0}{t} \implies vt = x - x_0 \implies
      \boxed{x = x_0 + vt}
    \]
    \item Reforce: \'e uma equa\c{c}\~ao de \textbf{reta} ($x$ em fun\c{c}\~ao de $t$).
  \end{itemize}

  \medskip
  \textbf{2. Gr\'aficos} (5--7 min)
  \begin{itemize}[nosep]
    \item Desenhe no quadro (use cor diferente para cada curva):
      \begin{itemize}[nosep]
        \item $x \times t$: reta com inclina\c{c}\~ao $v$ (positiva ou negativa).
        \item $v \times t$: reta horizontal.
      \end{itemize}
    \item Peça \`a turma: ``o que \'e a \'area sob o gr\'afico $v \times t$?''
          Guie at\'e chegarem em: \'area $=$ deslocamento.
  \end{itemize}

  \medskip
  \textbf{3. Exemplo: encontro de dois m\'oveis} (10--12 min)
  \begin{itemize}[nosep]
    \item Enunciado (dite ou projete):
          Dois carros partem de cidades A e B (distância 240 km)
          em sentidos opostos, com $v_1 = 80$ km/h e $v_2 = 40$ km/h.
          Quando e onde se encontram?
    \item Construa no quadro com os alunos:
    \[
      x_1 = 80t \qquad x_2 = 240 - 40t
    \]
    \[
      x_1 = x_2 \implies 120t = 240 \implies t = 2 \text{ h}
      \implies x = 160 \text{ km de A}
    \]
    \item Destaque: ``montar as equa\c{c}\~oes hor\'arias \'e o primeiro
          passo em qualquer problema de encontro.''
  \end{itemize}
\end{tcolorbox}

\begin{tcolorbox}[projetorbox]
  \textbf{Se o quadro ficar poluído com os gráficos:}
  projete a p\'agina de gr\'aficos do MRU das notas de aula.
  Aponte para cada curva enquanto explica o significado.

  \medskip
  \textbf{Alternativa:} manter o projetor desligado neste bloco e
  trabalhar s\'o no quadro --- o MRU \'e simples o suficiente para isso.
\end{tcolorbox}

\begin{tcolorbox}[perguntabox]
  Ap\'os o exemplo de encontro:
  \begin{itemize}[nosep]
    \item ``O que mudaria se os carros partissem no mesmo sentido?''
    \item ``Qual \'e a inclina\c{c}\~ao do gr\'afico $x \times t$ de um
          carro parado?''
  \end{itemize}
  \textit{Respostas: equa\c{c}\~ao de afastamento, n\~ao de encontro;
  inclina\c{c}\~ao zero ($v = 0$).}
\end{tcolorbox}

\begin{tcolorbox}[dicabox]
  No ENEM, quest\~oes de MRU frequentemente vêm na forma de gr\'aficos
  $v \times t$ ou $x \times t$ para interpretar --- n\~ao como conta direta.
  Treine a leitura gr\'afica tanto quanto a equa\c{c}\~ao.
\end{tcolorbox}

\begin{tcolorbox}[cortebox]
  Corte o exemplo de encontro se precisar --- ele reaparece nas
  quest\~oes do Bloco 6. Mantenha pelo menos a dedu\c{c}\~ao e os gr\'aficos.
\end{tcolorbox}

\newpage

% ===================================================================
\section*{Bloco 3 \hfill\colorbox{verdeok!30}{\small ESSENCIAL}}
\addcontentsline{toc}{section}{Bloco 3 --- MRUV}
\section*{\hspace{1.2em}MRUV --- Movimento Retil\'ineo Uniformemente Variado}
% ===================================================================
\tempo{40--50}

\begin{tcolorbox}[blocobox={Objetivos}]
  Deduzir as tr\^es equa\c{c}\~oes do MRUV usando \'algebra pura
  (sem c\'alculo). Interpretar gr\'aficos. Resolver problema de frenagem.
  Este \'e o bloco mais denso --- planeje o tempo com cuidado.
\end{tcolorbox}

\begin{tcolorbox}[quadrobox]
  \textbf{1. Defini\c{c}\~ao e 1\textordfeminine{} equa\c{c}\~ao} (5 min)
  \begin{itemize}[nosep]
    \item ``No MRUV, $a = $ cte $\neq 0$.''
    \item Dedução ao vivo:
    \[
      a = \frac{v - v_0}{t} \implies \boxed{v = v_0 + at}
    \]
  \end{itemize}

  \medskip
  \textbf{2. 2\textordfeminine{} equa\c{c}\~ao: via \'area do trap\'ezio}
  (10--12 min) \textit{[momento central da aula]}
  \begin{enumerate}[nosep]
    \item Desenhe o gr\'afico $v \times t$ no quadro com dois eixos e
          uma reta de $v_0$ at\'e $v$ (cor viva).
    \item Preencha a \'area sob a reta e diga: ``essa \'area \'e o
          deslocamento --- vamos calcular.''
    \item A figura \'e um trap\'ezio de bases $v_0$ e $v$ e altura $t$:
    \[
      \Delta x = \frac{v_0 + v}{2} \cdot t
    \]
    \item Substitua $v = v_0 + at$:
    \[
      \Delta x = \frac{v_0 + v_0 + at}{2}\cdot t = v_0 t + \tfrac{1}{2}at^2
    \]
    \[
      \boxed{x = x_0 + v_0 t + \tfrac{1}{2}at^2}
    \]
    \item Reforce: ``n\~ao memorizamos --- deduzimos. E a dedu\c{c}\~ao
          usa apenas \'area de trap\'ezio.''
  \end{enumerate}

  \medskip
  \textbf{3. 3\textordfeminine{} equa\c{c}\~ao: Torricelli} (7--8 min)
  \begin{enumerate}[nosep]
    \item ``Queremos $v$ em fun\c{c}\~ao de $\Delta x$, sem $t$.''
    \item Da 1\textordfeminine{}: $t = (v - v_0)/a$.
    \item Substitua em $\Delta x = \tfrac{v_0+v}{2}\cdot t$:
    \[
      \Delta x = \frac{v_0+v}{2}\cdot\frac{v-v_0}{a}
               = \frac{v^2-v_0^2}{2a}
    \]
    \[
      \boxed{v^2 = v_0^2 + 2a\,\Delta x}
    \]
  \end{enumerate}

  \medskip
  \textbf{4. Gr\'aficos do MRUV} (5 min)
  \begin{itemize}[nosep]
    \item Desenhe rapidamente os 3 gr\'aficos lado a lado:
          $x\times t$ (par\'abola), $v\times t$ (reta inclinada),
          $a\times t$ (constante).
    \item Inclina\c{c}\~ao do $v\times t$ \'e $a$.
    \item Raiz da par\'abola no $x\times t$ \'e quando o m\'ovel para.
  \end{itemize}

  \medskip
  \textbf{5. Exemplo: frenagem} (7--8 min)
  \begin{itemize}[nosep]
    \item Enunciado: carro a 72 km/h freia com $a = -5$ m/s\textsuperscript{2}.
          Qual a distância até parar?
    \item Convers\~ao ao vivo: 72 km/h $= 20$ m/s.
    \item Torricelli ($v = 0$):
    \[
      0 = 400 + 2(-5)\Delta x \implies \Delta x = 40\text{ m}
    \]
    \item Coment\'ario: ``esse \'e o tipo de problema que aparece no
          ENEM com contexto de seguran\c{c}a no tr\^ansito.''
  \end{itemize}
\end{tcolorbox}

\begin{tcolorbox}[projetorbox]
  \textbf{Momentos-chave para usar o projetor neste bloco:}
  \begin{enumerate}[nosep]
    \item \textbf{O trap\'ezio no $v \times t$:} se o desenho no quadro
          n\~ao ficar claro, projete o gr\'afico da p\'agina de MRUV
          das notas de aula.
          Aponte para as bases $v_0$ e $v$, a altura $t$, e a \'area
          destacada.
    \item \textbf{Os 3 gr\'aficos lado a lado:} projete-os ap\'os
          desenhá-los no quadro, para os alunos conferirem.
  \end{enumerate}
\end{tcolorbox}

\begin{tcolorbox}[perguntabox]
  Ap\'os deduzir a 1\textordfeminine{} equa\c{c}\~ao:
  \begin{itemize}[nosep]
    \item ``Se $a > 0$ e $v_0 > 0$, o que acontece com a velocidade?
          E se $a < 0$ e $v_0 > 0$?''
  \end{itemize}
  Antes de Torricelli:
  \begin{itemize}[nosep]
    \item ``Por que \'e \'util ter uma equa\c{c}\~ao sem o tempo?
          Em que situa\c{c}\~ao voc\^es n\~ao sabem o tempo?''
  \end{itemize}
  Ap\'os a frenagem:
  \begin{itemize}[nosep]
    \item ``Se a velocidade dobrar (144 km/h), a dist\^ancia de
          frenagem dobra tamb\'em?''
          \textit{N\~ao! Quadruplica. Deixe os alunos calcularem.}
  \end{itemize}
\end{tcolorbox}

\begin{tcolorbox}[dicabox]
  \textbf{A dedu\c{c}\~ao pela \'area \'e o coração deste bloco.}
  N\~ao pule para a f\'ormula diretamente. Quando os alunos entendem
  que $\Delta x$ \'e uma \'area, eles nunca mais confundem as equa\c{c}\~oes
  do MRUV.

  \medskip
  Se a turma parecer perdida no trap\'ezio: fa\c{c}a primeiro para
  velocidade constante (ret\^angulo) e depois generalize para MRUV
  (trap\'ezio). A progress\~ao ajuda.
\end{tcolorbox}

\begin{tcolorbox}[cortebox]
  Se precisar cortar: elimine o exemplo de frenagem deste bloco ---
  ele aparece como Quest\~ao 4 no Bloco 6.
  Mantenha \textbf{todas} as tr\^es dedu\c{c}\~oes.
\end{tcolorbox}

\newpage

% ===================================================================
\section*{Bloco 4 \hfill\colorbox{verdeok!30}{\small ESSENCIAL}}
\addcontentsline{toc}{section}{Bloco 4 --- Queda Livre e Lançamento Vertical}
\section*{\hspace{1.2em}Queda Livre e Lan\c{c}amento Vertical}
% ===================================================================
\tempo{20--25}

\begin{tcolorbox}[blocobox={Objetivos}]
  Mostrar que queda livre \'e MRUV com $a = g$.
  Deduzir as equa\c{c}\~oes e a simetria do lan\c{c}amento vertical.
\end{tcolorbox}

\begin{tcolorbox}[quadrobox]
  \textbf{1. Galileu e a queda livre} (3--4 min)
  \begin{itemize}[nosep]
    \item Conte brevemente: Galileu derrubou objetos de massas diferentes
          da Torre de Pisa e chegou \`a conclus\~ao de que todos caem com
          a mesma acelera\c{c}\~ao (na aus\^encia do ar).
    \item Escreva: $a = g \approx 10$ m/s\textsuperscript{2}, para baixo.
    \item ``\'E apenas MRUV com $v_0 = 0$ e $a = g$.''
    \item Escreva as equa\c{c}\~oes simplificadas (eixo positivo para baixo):
    \[
      v = gt \qquad h = \tfrac{1}{2}gt^2 \qquad v^2 = 2gh
    \]
  \end{itemize}

  \medskip
  \textbf{2. Exemplo: penhasco} (5 min)
  \begin{itemize}[nosep]
    \item Pedra cai de 80 m. Calcule tempo e velocidade de impacto.
    \[
      80 = 5t^2 \implies t = 4\text{ s} \qquad v = 10\times4 = 40\text{ m/s}
    \]
  \end{itemize}

  \medskip
  \textbf{3. Lan\c{c}amento vertical para cima} (8--10 min)
  \begin{itemize}[nosep]
    \item ``E se jogarmos para cima com $v_0 > 0$? Gravidade ainda \'e para
          baixo.'' Eixo positivo agora para cima $\to a = -g$.
    \item Escreva as equa\c{c}\~oes e destaque no quadro:
    \[
      v = v_0 - gt \qquad y = v_0 t - \tfrac{1}{2}gt^2
    \]
    \item Ponto mais alto ($v = 0$):
    \[
      t_{\text{sub}} = \frac{v_0}{g} \qquad H = \frac{v_0^2}{2g}
    \]
    \item Desenhe a trajet\'oria no quadro: seta para cima, desacelera,
          para no topo, acelera de volta.
    \item Simetria: $t_{\text{sub}} = t_{\text{des}}$; velocidade de
          retorno tem mesmo m\'odulo.
    \item Exemplo: bola com $v_0 = 30$ m/s $\to H = 45$ m, $T = 6$ s.
  \end{itemize}
\end{tcolorbox}

\begin{tcolorbox}[projetorbox]
  \textbf{Opcional:} projete o gr\'afico $v \times t$ do lan\c{c}amento
  vertical (reta de inclina\c{c}\~ao $-g$, cruzando zero no ponto mais
  alto e ficando negativa na descida). Isso ajuda os alunos a visualizarem
  a simetria.
\end{tcolorbox}

\begin{tcolorbox}[perguntabox]
  Antes do lan\c{c}amento para cima:
  \begin{itemize}[nosep]
    \item ``Se jogarmos uma bola para cima, ela fica quanto tempo no ar?
          Depende da massa?''
  \end{itemize}
  Ap\'os a simetria:
  \begin{itemize}[nosep]
    \item ``Se a bola passa pela posi\c{c}\~ao $y = 20$ m subindo, com
          que velocidade ela passa por $y = 20$ m descendo?''
          \textit{Mesmo m\'odulo, sinal oposto.}
  \end{itemize}
\end{tcolorbox}

\begin{tcolorbox}[dicabox]
  Muitos alunos acham que no ponto mais alto o objeto ``para'' e a
  gravidade ``para de agir''. Reforce: $v = 0$ mas $a = g \neq 0$.
  Um exemplo bom: o carro ao frear --- diminui a velocidade mas a
  acelera\c{c}\~ao (desacelera\c{c}\~ao) n\~ao \'e zero.
\end{tcolorbox}

\newpage

% ===================================================================
\section*{Bloco 5 \hfill\colorbox{laranjaalerta!40}{\small AVAN\c{C}ADO}}
\addcontentsline{toc}{section}{Bloco 5 --- Lançamento Oblíquo}
\section*{\hspace{1.2em}Lan\c{c}amento Obl\'iquo}
% ===================================================================
\tempo{35--40}

\begin{tcolorbox}[blocobox={Objetivos}]
  Decompor o movimento em horizontal (MRU) e vertical (MRUV).
  Deduzir tempo de voo, alcance e altura m\'axima.
  Este bloco usa mais o projetor --- a decomposi\c{c}\~ao vetorial
  \'e dif\'icil de ilustrar bem s\'o no quadro.
\end{tcolorbox}

\begin{tcolorbox}[projetorbox]
  \textbf{Use o projetor desde o in\'icio deste bloco.}

  \begin{enumerate}[nosep]
    \item \textbf{Figura da decomposi\c{c}\~ao de $v_0$:} mostre o vetor
          $v_0$ com \^angulo $\theta$, e as componentes $v_{0x} = v_0\cos\theta$
          e $v_{0y} = v_0\operatorname{sen}\theta$.
          Use a figura das notas de aula ou desenhe no quadro enquanto
          projeta.
    \item \textbf{Trajet\'oria parab\'olica:} projete a curva com as
          setas de $v_x$ (constante) e $v_y$ (que diminui e inverte).
          Isso \'e cr\'itico para a intui\c{c}\~ao.
    \item \textbf{Ponto mais alto:} destaque que $v_y = 0$ (mas $v_x$
          continua!) e a velocidade n\~ao \'e zero --- apenas vertical.
  \end{enumerate}
\end{tcolorbox}

\begin{tcolorbox}[quadrobox]
  \textbf{1. Princ\'ipio da independ\^encia} (5 min)
  \begin{itemize}[nosep]
    \item ``O movimento horizontal e o vertical n\~ao interferem um no
          outro.''
    \item Escreva lado a lado no quadro:

    \medskip
    \begin{tabular}{ll}
      \textbf{Horizontal (MRU)} & \textbf{Vertical (MRUV)} \\
      $v_x = v_0\cos\theta$ & $v_y = v_0\operatorname{sen}\theta - gt$ \\
      $x = v_0\cos\theta\cdot t$ &
        $y = v_0\operatorname{sen}\theta\cdot t - \tfrac{1}{2}gt^2$ \\
    \end{tabular}
  \end{itemize}

  \medskip
  \textbf{2. Tempo de voo $T$} (7 min)
  \begin{itemize}[nosep]
    \item ``O proj\'etil cai quando $y = 0$ de novo.'' Fatore:
    \[
      0 = v_0\operatorname{sen}\theta\cdot T
          - \tfrac{1}{2}g T^2
        = T\!\left(v_0\operatorname{sen}\theta - \tfrac{g}{2}T\right)
    \]
    \item Como $T \neq 0$: $T = 2v_0\operatorname{sen}\theta/g$.
  \end{itemize}

  \medskip
  \textbf{3. Alcance $R$} (5 min)
  \[
    R = v_0\cos\theta\cdot T = \frac{v_0^2\cdot 2\operatorname{sen}\theta\cos\theta}{g}
      = \frac{v_0^2\operatorname{sen}(2\theta)}{g}
  \]
  Conclua: m\'aximo em $\theta = 45^\circ$.

  \medskip
  \textbf{4. Altura m\'axima $H$} (5 min)
  \begin{itemize}[nosep]
    \item $v_y = 0 \implies t_H = v_0\operatorname{sen}\theta/g$.
    \item Torricelli na vertical:
    $H = (v_0\operatorname{sen}\theta)^2/(2g)$.
  \end{itemize}

  \medskip
  \textbf{5. Exemplo} (8--10 min)
  \begin{itemize}[nosep]
    \item Bola chutada: $v_0 = 20$ m/s, $\theta = 30^\circ$.
          Calcule $H$, $R$ e $T$ ao vivo com a turma.
  \end{itemize}
\end{tcolorbox}

\begin{tcolorbox}[perguntabox]
  Antes de calcular o alcance:
  \begin{itemize}[nosep]
    \item ``Para qual \^angulo o alcance \'e maior? Por qu\^e?''
    \item ``O que acontece com o alcance se dobrarmos $v_0$?''
          \textit{Quadruplica: $R \propto v_0^2$.}
  \end{itemize}
  Ap\'os calcular $R$ para $30^\circ$:
  \begin{itemize}[nosep]
    \item ``Para $60^\circ$ (complementar de $30^\circ$), o alcance
          muda?'' \textit{N\~ao --- mostre pelo $\operatorname{sen}(2\times60^\circ)
          = \operatorname{sen}(120^\circ) = \operatorname{sen}(60^\circ)$.}
  \end{itemize}
\end{tcolorbox}

\begin{tcolorbox}[dicabox]
  O ponto mais dif\'icil aqui \'e convencer os alunos de que
  horizontal e vertical s\~ao \textbf{realmente independentes}.
  Uma boa imagem: imagine dois caminh\~oes na mesma estrada, um
  deles acelerando (vertical) e o outro na velocidade constante
  (horizontal) --- nenhum interfere no outro.
\end{tcolorbox}

\begin{tcolorbox}[cortebox]
  Se precisar cortar: elimine a dedu\c{c}\~ao de $H$ (fique s\'o com
  $T$ e $R$) e reduza o exemplo para s\'o $R$ e $T$.
  O lan\c{c}amento obl\'iquo \'e menos frequente no ENEM do que
  queda livre e MRUV.
\end{tcolorbox}

\newpage

% ===================================================================
\section*{Bloco 6 \hfill\colorbox{verdeok!30}{\small ESSENCIAL}}
\addcontentsline{toc}{section}{Bloco 6 --- Fixação com Questões ENEM}
\section*{\hspace{1.2em}Fixa\c{c}\~ao com Quest\~oes de ENEM}
% ===================================================================
\tempo{20--25}

\begin{tcolorbox}[blocobox={Objetivos}]
  Conectar o conte\'udo ao formato ENEM. Mostrar que as quest\~oes
  pedem leitura e interpreta\c{c}\~ao, n\~ao apenas aplica\c{c}\~ao direta.
\end{tcolorbox}

\begin{tcolorbox}[projetorbox]
  Projete as quest\~oes uma de cada vez a partir do arquivo
  \texttt{notas\_aula\_cinematica.pdf}, Se\c{c}\~ao 7 (\'ultimas p\'aginas).

  \medskip
  \textbf{Din\^amica sugerida para cada quest\~ao:}
  \begin{enumerate}[nosep, label=\arabic*.]
    \item Projete o enunciado. D\^e 2--3 minutos para os alunos tentarem.
    \item Pergunte quem marcou qual alternativa.
    \item Resolva \textbf{no quadro} (n\~ao apenas leia a resposta).
    \item Aponte qual habilidade/equa\c{c}\~ao aquela quest\~ao testou.
  \end{enumerate}
\end{tcolorbox}

\begin{tcolorbox}[quadrobox]
  \textbf{Quest\~oes sugeridas para este bloco} (selecione conforme o tempo):

  \medskip
  \begin{tabular}{lll}
    \toprule
    \# & T\'opico & Tempo estimado \\
    \midrule
    Q2 & Leitura de gr\'afico $v\times t$ & 5 min \\
    Q3 & MRUV + Torricelli & 5 min \\
    Q5 & Queda livre & 4 min \\
    Q6 & MRUV: dist. de frenagem & 4 min \\
    \bottomrule
  \end{tabular}

  \medskip
  Resolva cada uma no quadro, passo a passo, como faria numa prova.
\end{tcolorbox}

\begin{tcolorbox}[dicabox]
  \textbf{N\~ao resolva todas de uma vez.} Prefira resolver 3 quest\~oes
  com discuss\~ao a correr por 8 sem conversa. O erro mais comum no
  ENEM de f\'isica \'e leitura errada do enunciado, n\~ao erro de conta
  --- treine isso.
\end{tcolorbox}

\begin{tcolorbox}[cortebox]
  Se sobrou menos de 10 min: fa\c{c}a apenas Q2 (gr\'afico) --- \'e a
  habilidade mais testada no ENEM e a mais descuidada nas revis\~oes.
\end{tcolorbox}

\newpage

% ===================================================================
\section*{Encerramento e Pr\'oximos Passos}
\addcontentsline{toc}{section}{Encerramento}
% ===================================================================
\tempo{5}

\begin{tcolorbox}[quadrobox]
  Deixe no quadro (ou projete) o resumo das f\'ormulas da Se\c{c}\~ao 6
  das notas de aula. Aponte cada equa\c{c}\~ao e reforce:
  \begin{itemize}[nosep]
    \item ``Voc\^es sabem de onde todas essas equa\c{c}\~oes v\^em.''
    \item ``N\~ao precisam memorizar Torricelli --- s\'o precisam saber
          o que \'e \'area de trap\'ezio.''
  \end{itemize}
\end{tcolorbox}

\begin{tcolorbox}[perguntabox]
  Encerramento r\'apido com a turma:
  \begin{itemize}[nosep]
    \item ``Qual equa\c{c}\~ao voc\^es acham que vai aparecer mais no ENEM?''
    \item ``Qual parte ainda ficou com d\'uvida?''
  \end{itemize}
\end{tcolorbox}

\begin{tcolorbox}[dicabox]
  \textbf{Tarefa sugerida (se houver):} resolver as Quest\~oes 1, 7 e 8
  das notas de aula (MRU, lan\c{c}. obl\'iquo e \^angulos complementares)
  que n\~ao foram vistas em sala.
\end{tcolorbox}

\newpage

% ===================================================================
% GUIA RÁPIDO DE ADAPTAÇÃO
% ===================================================================
\section*{Guia R\'apido de Adapta\c{c}\~ao por Tempo Dispon\'ivel}
\addcontentsline{toc}{section}{Guia de Adaptação por Tempo}

\renewcommand{\arraystretch}{1.6}
\begin{tabularx}{\linewidth}{lX}
  \toprule
  \textbf{Tempo dispon\'ivel} & \textbf{Blocos a incluir} \\
  \midrule
  $\approx 1$h &
    B0 (10 min) + B1 (20 min) + B3 s\'o equa\c{c}\~oes (15 min)
    + 2 quest\~oes do B6 (15 min) \\
  $\approx 1$h30 &
    B0 + B1 + B2 (compacto) + B3 + 2 quest\~oes do B6 \\
  $\approx 2$h &
    B0 + B1 + B2 + B3 + B4 + 3 quest\~oes do B6 \\
  $\approx 2$h30 &
    Todos os blocos, com ritmo normal \\
  $\approx 3$h &
    Todos os blocos, com tempo para perguntas e mais exerc\'icios \\
  \bottomrule
\end{tabularx}

\bigskip

\begin{tcolorbox}[colback=azulphysics!5!white, colframe=azulphysics, arc=3mm,
  fonttitle=\bfseries, title={Prioridade dos t\'opicos para ENEM}]
  \begin{enumerate}[nosep]
    \item \textbf{Leitura de gr\'aficos} $v \times t$ e $x \times t$
          (muito frequente)
    \item \textbf{MRUV} --- Torricelli e dist\^ancia de frenagem
    \item \textbf{Queda livre} e lan\c{c}amento vertical
    \item \textbf{MRU} --- encontro de m\'oveis e velocidade relativa
    \item \textbf{Lan\c{c}amento obl\'iquo} (menos frequente, mas cai)
  \end{enumerate}
\end{tcolorbox}

\end{document}
